\section{The same obstruction in the printed twisted Hilbert realization}
\label{sec:printed-projection-obstruction}

\paragraph{Status and purpose.}
The complete author reading and two independent whole-proof
readings have passed; their actual check scopes are recorded separately.  MA-PI-001 fixes the corrected centered
symmetric norm.  Here the original printed Hilbert realization is
derived from its own definitions, and its orthogonal projector is
tested directly.  The change of realization below is an actual
isometry, not an appeal to equivalence of norms.

Retain $\chi=\chi_4$, $S=\{\infty,2\}$, the actual finite vector
$\theta_2$, the canonical adelic source $\varphi_f$, and the
physical cutoff $\mathfrak B_{S,Y}$ from (MA-PI.1--9).
The outside-$S$ factors of $\varphi_f$ are
$\mathbf1_{\mathbb Z_p}$.  For every odd Schwartz $f$, both
adelic polar conditions hold at the ramified finite place.
Let $A$ denote the source's twist parameter, to distinguish it
from the symmetric width $a$ in MA-PI-001.

\paragraph{The exact source coordinates and metric.}
Masullo Definition~2.2.3.1, equation~(2.2.3.1),
printed/PDF pp.~81--82, defines the \emph{uncentered} theta
function
\[
 \Phi_{\mathbb Q}(\varphi)(g)
       =\sum_{\xi\in\mathbb Q^\times}\varphi(\xi g).
 \tag{MA-PI.32}
\]
Definition~9.1.2.1, equation~(9.1.2.3), p.~399, defines
the norm on $W^{s,2}(C_{\mathbb Q})$ by
\[
 \|F\|_{s,2}^2
     =\int_{C_{\mathbb Q}}|F(g)|^2
              (1+|\log|g||^2)^{s/2}\,d\mu(g).
 \tag{MA-PI.33}
\]
The compact-character component is denoted
$\mathcal H_s^\chi=W^{s,2}(C_{\mathbb Q})(\chi|_{C_{\mathbb Q}^1})$.
Definition~10.1.1.1 and its preceding equations
(10.1.1.1--3), pp.~434--435, realize the twist by
the actual function $|g|^A\Phi_{\mathbb Q}(\varphi)(g)$.
The same twist appears in the invariant PW level of
Definition~11.2.1.3, equations~(11.2.1.7--8), pp.~488--489.
Let $Q^{\mathrm{src}}_{S,Y,A,s}$ be that printed
orthogonal projection in $\mathcal H_s^\chi$.

Choose the positive archimedean norm splitting $\nu$ and write
$g=\nu(e^x)k$, where $k\in C_{\mathbb Q}^1$.
As in the paragraph following (9.1.2.3), the Haar measure
has the exact product decomposition $dx\,d\mu_1(k)$.
Put $m=\mu_1(C_{\mathbb Q}^1)>0$; this retains the source's
Haar normalization rather than silently assigning it mass one.
For the actual canonical character source one has
\[
 \begin{split}
 (T_Af)(\nu(e^x)k)
   &:=|g|^A\Phi_{\mathbb Q}(\varphi_f)(g)\\
   &=\chi(k)^{-1}e^{b x}Y_f(x),
          \qquad b=A-\tfrac12,
 \end{split}
 \tag{MA-PI.34}
\]
where $Y_f$ is exactly (MA-PI.2).
Indeed the actual finite vector restricts the theta sum to
integers with coefficient $\chi_4(n)$, and the two signs
give the factor $2$ in (MA-PI.2).  Its factor $e^{x/2}$
is the centering, while the source twist is $e^{Ax}$.
The compact-character dependence is that of MA-CGP.2--7;
here $\chi^{-1}=\chi$, so no inverse-character ambiguity
affects (MA-PI.34).

Define the one-sided Hilbert space and an actual unitary map by
\[
 \begin{split}
 H^{\mathrm{one}}_{b,s}
    &=L^2\bigl(\mathbb R,e^{2bx}(1+x^2)^{s/2}\,dx\bigr),\\
 (U_bF)(\nu(e^x)k)
    &=m^{-1/2}\chi(k)^{-1}e^{bx}F(x).
 \end{split}
 \tag{MA-PI.35}
\]
Equations (MA-PI.33) and $|\chi(k)|=1$ prove that $U_b$
is isometric onto $\mathcal H_s^\chi$: the compact-character
projector shows that every function in that component has
the displayed form almost everywhere.  Its inverse is the
well-defined fiber coefficient
\[
 (U_b^{-1}H)(x)=m^{-1/2}e^{-bx}
        \int_{C_{\mathbb Q}^1}H(\nu(e^x)k)\chi(k)\,d\mu_1(k).
\]
This formula is understood for almost every $x$, using
Fubini and Cauchy--Schwarz on the compact fiber.  In particular
\[
 T_A f=\sqrt m\,U_bY_f,\qquad
 \langle T_Af,T_Ah\rangle_{s,2}
     =m\int_{\mathbb R}Y_f(x)\overline{Y_h(x)}
                           e^{2bx}(1+x^2)^{s/2}\,dx.
 \tag{MA-PI.36}
\]
Consequently $U_b^{-1}Q^{\mathrm{src}}_{S,Y,A,s}U_b$
is the orthogonal projector onto the actual inverse image
of the printed PW range.  This assertion follows from
unitarity; no assertion about orthogonal projectors under
a merely bounded norm equivalence is used.

\paragraph{MA-PI-002 (failure for the original printed projection).}
For every $0<A<1$ and every $s\geq0$, there is $Y_0>0$
such that, for every $Y\geq Y_0$ and every smooth physical
profile (MA-PI.4), there is an odd Schwartz function $f_Y$
flat in a neighborhood of zero for which
\[
 Q^{\mathrm{src}}_{S,Y,A,s}T_Af_Y
                     \ne T_A(\mathfrak b_Yf_Y).
 \tag{MA-PI.37}
\]
The input $f_Y\otimes\theta_2$ is in the original
coordinate-flat LF source of MA-GA.2.  Its physical
cutoff output is genuine $S$-local Schwartz data satisfying
both polar conditions, so the right side is the actual
printed global theta map on that output.
In particular (MA-PI.37) applies to $A=1/2$, the exact
twist specified in the proof of Theorem~12.4.1.1 on p.~515.
The original printed canonical intertwining is therefore
false, without a norm-realization qualification.
This does not assert the failure of every different
Gram-corrected factorization.

\begin{proof}
\emph{1. Original source membership and physical action.}
The finite Fourier calculation (MA-PI.9) shows that the
canonical data have shell $|\xi_2|_2=4$ and that
\[
 \mathfrak B_{S,Y}(f\otimes\theta_2)
       =(\mathfrak b_Yf)\otimes\theta_2,
 \quad
 \mathfrak b_Yf=\mathcal F_\infty^{-1}
                      \bigl(\rho_Y(4|\xi|)\widehat f(\xi)\bigr).
 \tag{MA-PI.38}
\]
This is a literal identity on the stated canonical Schwartz
data.  The norm multiplier is even, so the equivalent
expression $\mathcal F_\infty M_\rho\mathcal F_\infty^{-1}$
gives the same operator: $\mathcal F_\infty^2$ is reflection
and reflection commutes with $M_\rho$.  The same observation
applies to coordinatewise reflected finite Fourier conventions.
Thus the unchanged physical norm cutoff of the source is used.

If $\widehat h$ is compact smooth and supported in
$|\xi_\infty|<Y/4$, the entire $S$-Fourier support is
compact with product norm $<Y$.  Both polar conditions
hold at $2$.  Definition~7.3.2.3, (7.3.2.12--17),
pp.~354--355, with $S'=S$, therefore puts the actual
theta image in the source invariant PW level.  Twisting
by $|g|^A$ and taking the closure gives
$T_Ah\in\operatorname{Ran}Q^{\mathrm{src}}_{S,Y,A,s}$.
This is direct membership, not an assumed equality of
different possible PW ranges.

\emph{2. A nonzero cross pairing for this exact one-sided line.}
Put $b=A-1/2\in(-1/2,1/2)$.
Equations (MA-PI.10--20) already prove, for every such $b$,
the exact one-sided theta Gram identity and that its
inverse-Fourier distribution is not compactly supported.
Those formulas were derived for the whole open range of
$b$, before the specialization $b=-a$ in MA-PI-001.
Their distributional separation argument therefore
supplies odd Schwartz functions $f,h$ whose Fourier
transforms are compact smooth with respective supports
\[
 \begin{gathered}
 \operatorname{supp}\widehat f\subset\{|\xi|>1/2\},\qquad
 \operatorname{supp}\widehat h\subset\{0<|\xi|<1/8\},\\
 D_b:=\int_{\mathbb R}Y_f(x)\overline{Y_h(x)}e^{2bx}\,dx\ne0.
 \end{gathered}
 \tag{MA-PI.39}
\]
These are genuine Schwartz source vectors.  No zero-location
or mean-square hypothesis is used; the separation follows
from the proved factorial growth of the entire beta product.

\emph{3. Dilation in the printed polynomial norm.}
Use the actual dilations $f_L(t)=e^{L/2}f(e^Lt)$ and
$h_L(t)=e^{L/2}h(e^Lt)$, and set $Y=e^L$.
Their Fourier supports and the two profile plateaus give
\[
 \mathfrak b_Y f_L=0,\qquad
 \mathfrak b_Y h_L=h_L,\qquad
 T_Ah_L\in\operatorname{Ran}Q^{\mathrm{src}}_{S,Y,A,s}.
 \tag{MA-PI.40}
\]
Exact theta scaling gives $Y_{f_L}(x)=Y_f(x+L)$ and the
same formula for $h$.  Substituting into the actual
printed inner product (MA-PI.36), for $L\geq1$, yields
\[
 \begin{split}
 &m^{-1}e^{2bL}L^{-s}
        \langle T_A f_L,T_A h_L\rangle_{s,2}\\
 &\qquad=\int_{\mathbb R}Y_f(x)\overline{Y_h(x)}e^{2bx}
                  \frac{(1+(x-L)^2)^{s/2}}{L^s}\,dx
                         \longrightarrow D_b\ne0.
 \end{split}
 \tag{MA-PI.41}
\]
The polynomial ratio tends to one and is bounded by
$C_s(2+|x|)^s$.  Both theta images belong to $\mathcal E$,
so dominated convergence proves (MA-PI.41).  The pairing
is therefore nonzero for every sufficiently large $L$,
independently of the choice of transition profile.
Orthogonality of the actual printed projector and
(MA-PI.40) imply
\[
 \bigl\langle
 Q^{\mathrm{src}}_{S,Y,A,s}T_Af_L
       -T_A(\mathfrak b_Y f_L),\ T_Ah_L\bigr\rangle_{s,2}
        =\langle T_Af_L,T_Ah_L\rangle_{s,2}\ne0.
 \tag{MA-PI.42}
\]

\emph{4. A finite witness in the original flat source.}
Fix such $Y$ and its smooth profile, and choose the even
compact cutoff $\omega=1$ near zero as in (MA-PI.27).
Set $f_{L,\varepsilon}=f_L-r_\varepsilon$, where
$r_\varepsilon(t)=\omega(t/\varepsilon)f_L(t)
                  =\varepsilon g_\varepsilon(t/\varepsilon)$.
The $g_\varepsilon$ form a bounded family of odd compact
Schwartz functions.  Using (MA-PI.28) and the exact
one-sided norm in (MA-PI.36) gives
\[
 \|T_A r_\varepsilon\|_{s,2}
      \leq C\varepsilon^{3/2+b}
                    (1+|\log\varepsilon|)^{s/2}
       =C\varepsilon^{A+1}(1+|\log\varepsilon|)^{s/2}
                          \longrightarrow0.
 \tag{MA-PI.43}
\]
For completeness, the substitution $x=y+\log\varepsilon$
in the squared integral contributes $\varepsilon^3$
from theta scaling and $\varepsilon^{2b}$ from the
exponential weight.  The remaining polynomial ratio is
bounded by a constant times
$(1+|\log\varepsilon|)^s(1+|y|)^s$; uniform exponential
Schwartz bounds for $Y_{g_\varepsilon}$ finish the estimate.
The exponent $A+1$ is positive throughout the asserted range.

For the fixed profile, (MA-PI.29) proves that
$\mathfrak b_Y r_\varepsilon\to0$ in odd Schwartz space:
its Fourier transform is smooth with fixed compact support,
and every derivative tends uniformly to zero.  The actual
theta map is continuous into $\mathcal E$, and multiplication
by $e^{bx}$ is continuous from $\mathcal E$ into polynomial
$L^2$.  Hence
\[
       \|T_A(\mathfrak b_Y r_\varepsilon)\|_{s,2}\to0.
 \tag{MA-PI.44}
\]
The actual orthogonal projector is bounded.  Equations
(MA-PI.43--44) show that the pairing in (MA-PI.42), with
$f_L$ replaced by $f_{L,\varepsilon}$, converges to its
nonzero value.  Choose a sufficiently small positive
$\varepsilon$ and put $f_Y=f_{L,\varepsilon}$.
This function is Schwartz and flat near the archimedean
origin; its finite factor is supported away from zero.
It is therefore an actual vector of the original inherited
LF source, and it proves (MA-PI.37).  The right side
remains genuine Schwartz polar-free data by (MA-PI.38).
\end{proof}

\paragraph{Source and conclusion boundaries.}
The exact source definitions, compact Haar product decomposition,
twisting convention, invariant PW range and orthogonal projection
were read in full at the locators given above.  The physical
cutoff is the source's Step~II, p.~516, and composition on
p.~517, in Theorem~12.4.1.1.  Its proof explicitly uses
$W_{1/2}^{s,2}$ on p.~515.  Formula (MA-PI.34) therefore
identifies that actual case with $b=0$, without a centering
guess or an equivalent-norm substitution.

The finite Artin representation, inertia convention and
Grothendieck placement are exactly those recorded after
MA-PI-001; changing the analytic Hilbert realization does
not change that rank-one object.  The result determines
the specified canonical theta/physical-cutoff map identity
to be false in both the corrected and the literal printed
realizations.  It does not attribute an unprinted equation
to the source, and it does not rule out a different
Gram-corrected comparison, alter the valid nuclear returns,
or remove the arithmetic FT002 condition.
